% !TeX root = ../../tesis.tex
%%%%%%%%%%%%%%%%%%%%%%%%%%%%%%%%%%%%%%%%%%%%%%%%%%%%%%%%%%%%%%%%%%%
% SECCIÓN 5.3.4 — FACTOR DE DESVIACIÓN (DI)
%%%%%%%%%%%%%%%%%%%%%%%%%%%%%%%%%%%%%%%%%%%%%%%%%%%%%%%%%%%%%%%%%%%
\subsection{Factor de Desviación (\texorpdfstring{$DI$}{DI})}
\label{subsec:factor-desviacion}

Cuando alguien planea un viaje con el mapa en mano, la distancia en
línea recta entre origen y destino parece razonable. Lo que la red de
transporte impone en la práctica es otra cosa: rutas que obligan a
converger hacia el centro antes de salir hacia la periferia opuesta,
trayectos que multiplican el recorrido respecto al camino ideal. El
\textbf{Factor de Desviación} ($DI$) cuantifica exactamente esa
brecha: es la razón entre la distancia real recorrida sobre la
infraestructura —incluyendo la primera y última milla peatonal— y la
distancia teórica en línea recta entre el mismo par origen-destino.

\begin{equation}
    DI = \frac{D_{\text{red}} + D_{\text{caminata}}}
              {D_{\text{haversine}}(O, D)}
    \label{eq:factor-desviacion}
\end{equation}

\noindent donde:
\begin{itemize}
    \item $D_{\text{red}}$: distancia acumulada sobre el trazo real de
          la ruta óptima en tiempo.
    \item $D_{\text{caminata}}$: suma de la primera y última milla
          peatonal (acceso y egreso a la red).
    \item $D_{\text{haversine}}(O, D)$: distancia en línea recta entre
          origen $O$ y destino $D$, calculada con la fórmula de
          Haversine.
\end{itemize}

Un valor de $DI = 1.0$ indica una ruta perfectamente directa; valores
superiores a $1.5$ revelan ineficiencias topológicas significativas en
el par origen-destino analizado.

% ---------------------------------------------------------------
\subsubsection{Distribución en la Red}
\label{subsubsec:di-distribucion}
% ---------------------------------------------------------------

Sobre una muestra de \textbf{100 pares origen-destino} (semilla
reproducible $= 42$), el modelo arrojó los estadísticos globales del
Cuadro~\ref{tab:di-estadisticas}. La distribución resultante —con
media de $1.545$ y valores extremos entre $0.97$ y $4.34$— confirma
que el recorrido promedio en la red metropolitana supera en más del
50\,\% la distancia ideal.

\begin{table}[H]
    \centering
    \caption[Estadísticos globales del Factor de Desviación]{%
    Estadísticos descriptivos del $DI$ sobre 100 pares
    origen-destino de la red de la \zmvm{}.
    Corrida: 2026-04-27.}
    \label{tab:di-estadisticas}
    \begin{tabular}{|l|r|}
        \hline
        \textbf{Estadístico} & \textbf{Valor} \\
        \hline
        $n$ (rutas calculadas)     & 100   \\
        Media ($\bar{DI}$)         & 1.545 \\
        Desviación estándar        & 0.517 \\
        Mínimo                     & 0.97  \\
        Percentil 25               & 1.27  \\
        Mediana                    & 1.40  \\
        Percentil 75               & 1.67  \\
        Máximo                     & 4.34  \\
        \hline
    \end{tabular}
    \fuentefigura{Elaboración propia con \texttt{VFTModel}
    \citep{cabrera2026VFTModel}. Corrida: 2026-04-27.}
\end{table}

El registro completo de los 100 pares analizados —con sus distancias
individuales, valores de $DI$ y número de saltos— se incluye en el
Cuadro~\ref{tab:anx-di-corrida-completa} del
Anexo~\ref{anx:visualizaciones-analiticas}.

La Figura~\ref{fig:mapa-di} sitúa esa distribución sobre el
territorio, permitiendo leer la magnitud del rodeo como fenómeno
espacial antes de desagregarla por alcaldía.

\begin{figure}[H]
    \centering
    \includegraphics[width=0.90\textwidth]{%
        Figures/Mapas/Qgis_Resources/RedActual/Mapa_DF}
    \caption[Distribución espacial del Factor de Desviación ($DI$)
    por nodo]{%
    Factor de Desviación ($DI$) por nodo en la red de transporte
    masivo de la \zmvm{}. El gradiente verde-a-rojo señala las
    zonas donde la infraestructura impone mayores rodeos respecto
    a la distancia lineal ideal. Corrida: 2026-04-27.}
    \label{fig:mapa-di}
    \fuentefigura{Elaboración propia con \textbf{QGIS};
    datos: \texttt{VFTModel} \citep{cabrera2026VFTModel},
    GTFS \citep{adip_gtfs_2024} e INEGI 2023.}
\end{figure}

La Figura~\ref{fig:grafica-di-alcaldia} desagrega esa imagen
territorial mostrando el $DI$ promedio por alcaldía, evidenciando
cuáles demarcaciones concentran la mayor ineficiencia de trayecto.

\begin{figure}[H]
    \centering
    \includegraphics[width=0.92\textwidth]{%
        Figures/Mapas/Tableau_Resources/RedActual/VFTModel/UniGrafica/%
        Grafica_VFTModel_2026_Factor_desviacion}
    \caption[Factor de Desviación promedio por alcaldía de la
    \zmvm{}]{%
    Factor de Desviación ($DI$) promedio por alcaldía de la
    \zmvm{}, calculado sobre la muestra de 100 pares
    origen-destino. Las alcaldías periféricas del oriente y sur
    presentan los valores más altos, reflejando la ausencia de
    conectividad transversal en esos cuadrantes.
    Corrida: 2026-04-27.}
    \label{fig:grafica-di-alcaldia}
    \fuentefigura{Elaboración propia con \textbf{Tableau};
    datos: \texttt{VFTModel} \citep{cabrera2026VFTModel}.}
\end{figure}

El patrón es consistente con los indicadores anteriores: las
demarcaciones que registran menor cobertura y menor densidad de
nodos de alta Fuerza Capilar son también las que imponen los mayores
rodeos a sus usuarios.

% ---------------------------------------------------------------
\subsubsection{Casos de Estudio}
\label{subsubsec:di-casos}
% ---------------------------------------------------------------

Para anclar los estadísticos en viajes concretos, se seleccionaron
doce pares origen-destino estratégicos concentrados en las periferias
sur y oriente de la cuenca. El Cuadro~\ref{tab:di-casos} consolida
los valores de $DI$ para cada caso.

\begin{table}[H]
    \centering
    \caption[Factor de Desviación — doce casos de estudio]{%
    Factor de Desviación ($DI$) para doce pares origen-destino
    estratégicos de la \zmvm{}, seleccionados en las periferias
    sur y oriente. Corrida: 2026-04-27.}
    \label{tab:di-casos}
    \small
    \begin{tabularx}{\textwidth}{|l|X|X|r|r|r|}
        \hline
        \textbf{\#} & \textbf{Origen} & \textbf{Destino} &
        \textbf{$D_{\text{red}}$ (km)} &
        \textbf{$D_{\text{Hav.}}$ (km)} &
        \textbf{$DI$} \\
        \hline
         1 & Milpa Alta        & Autódromo / Centro-Oriente   & 29.18 & 22.59 & 1.29 \\
         2 & Milpa Alta        & Paseo de la Reforma (CBD)    & 37.73 & 27.50 & 1.37 \\
         3 & Milpa Alta        & Santa Fe (CBD)               & 36.42 & 28.62 & 1.27 \\
         4 & Milpa Alta        & FES Acatlán                  & 48.34 & 35.73 & 1.35 \\
         5 & Milpa Alta        & Ciudad Universitaria         & 26.25 & 19.96 & 1.32 \\
         6 & Milpa Alta        & Coyoacán                     & 25.24 & 18.86 & 1.34 \\
         7 & Santa Catarina    & Ciudad Universitaria         & 42.06 & 23.10 & 1.82 \\
         8 & Santa Catarina    & FES Acatlán                  & 45.64 & 32.16 & 1.42 \\
         9 & Milpa Alta        & Santa Catarina (Sur-Oriente) & 37.22 & 14.79 & 2.52 \\
        10 & Milpa Alta        & Ajusco                       & 29.19 & 20.56 & 1.42 \\
        11 & Ajusco            & Preparatoria 1               & 15.91 & 11.16 & 1.43 \\
        12 & Parque Los Olivos & Centro de Tláhuac            &  2.20 &  2.08 & 1.06 \\
        \hline
        \multicolumn{5}{|l|}{\textbf{Mejor caso (muestra aleatoria):} Plutarco Elías Calles $\to$ Fuente de Loreto} & \textbf{0.97} \\
        \multicolumn{5}{|l|}{\textbf{Peor caso (muestra aleatoria):} Carlos A. Vidal $\to$ 20 de Noviembre} & \textbf{4.34} \\
        \hline
    \end{tabularx}
    \fuentefigura{Elaboración propia con \texttt{VFTModel}
    \citep{cabrera2026VFTModel}. Corrida: 2026-04-27.}
\end{table}

Entre los doce casos, el Caso~9 registra el $DI$ más alto del
conjunto seleccionado ($DI = 2.52$): para ir de Milpa Alta a Santa
Catarina Yecahuizotl —dos puntos del sur-oriente separados por apenas
14.8~km en línea recta— la red obliga a recorrer 37.2~km, más del
doble de la distancia ideal. Ese valor no es un accidente de la
muestra; es la expresión métrica de la ausencia de conectividad
transversal en ese cuadrante de la cuenca.

% ---------------------------------------------------------------
\subsubsection{Vista Integral de la Red}
\label{subsubsec:di-vista-integral}
% ---------------------------------------------------------------

¿Cómo se distribuye el costo del rodeo sobre el territorio? La
Figura~\ref{fig:dashboard-di} responde esa pregunta articulando en
un solo panel la magnitud del $DI$ por alcaldía, la dispersión
estadística de la muestra y la huella geográfica de las rutas
analizadas.

\begin{figure}[H]
    \centering
    \includegraphics[width=0.95\textwidth]{%
        Figures/Mapas/Tableau_Resources/RedActual/VFTModel/Dashboards/%
        Dashboard_VFTModel_2026_Factor_desviacion}
    \caption[Dashboard del Factor de Desviación — red multimodal
    \zmvm{}]{%
    Vista integral del Factor de Desviación ($DI$) sobre la red de
    transporte masivo de la \zmvm{}, generada con el
    \termino{dashboard} analítico del \texttt{VFTModel}.
    Corrida: 2026-04-27.}
    \label{fig:dashboard-di}
    \fuentefigura{Elaboración propia con \textbf{Tableau};
    datos: \texttt{VFTModel} \citep{cabrera2026VFTModel}.}
\end{figure}

El análisis del \texttt{VFTModel} identifica en Tláhuac la alcaldía
donde ese costo estructural es más pronunciado. El Caso~7 lo
cuantifica con precisión: de Santa Catarina Yecahuizotl a Ciudad
Universitaria hay 23~km en línea recta, pero la red exige 42~km de
recorrido real ($DI = 1.82$), porque ninguna ruta de los datos
procesados por \texttt{Apímetro} conecta esos dos puntos del sur sin
internar el trayecto hacia el centro de la ciudad. No es un caso
límite ni un atípico estadístico: es la geometría de una red sin
conectividad transversal en ese cuadrante, medida con los mismos
instrumentos que describen el resto del capítulo.

\begin{figure}[H]
    \centering
    \includegraphics[width=0.85\textwidth]{%
        Figures/Cap5/FactorDesviacion/Mapa/NB04_mapa10_caso07.png}
    \caption[Caso~7: Santa Catarina Yecahuizotl → Ciudad
    Universitaria ($DI = 1.82$)]{%
    Caso~7: ruta óptima de Santa Catarina Yecahuizotl a Ciudad
    Universitaria. La red obliga a un recorrido de 42.06~km para
    cubrir una distancia lineal de 23.10~km ($DI = 1.82$), al no
    existir conectividad transversal directa en el cuadrante
    sur-oriente de la \zmvm{}. Corrida: 2026-04-27.}
    \label{fig:di-caso07}
    \fuentefigura{Elaboración propia con \texttt{VFTModel}
    \citep{cabrera2026VFTModel}.}
\end{figure}

Los cuatro indicadores de Fase~1 —impedancia, cobertura, Fuerza
Capilar y Factor de Desviación— describen colectivamente una red
metropolitana con capacidades desiguales: bien articulada en el
centro, con déficit progresivo hacia la periferia sur y oriente, y
sin conectividad transversal en ese cuadrante. Esa descripción es
topológica: se basa en la estructura del grafo, sin considerar las
condiciones reales de operación. La Fase~2 incorpora esas condiciones
mediante dos parámetros de ponderación dinámica —el Coeficiente de
Fricción Vial ($CF$) y la penalización por transferencia ($T_b$)—
que ajustan el modelo a la experiencia de viaje efectiva.
